\chapter{The Ethical Horizon: AI for Global Climate Resilience}

The technical journey of this volume—from the fundamental physics of the atmospheric engine to the exascale architectures of planetary foundation models—culminates in a final, inescapable frontier: the \textbf{Ethical Horizon}. 

\section{The Energy-Climate Paradox: Carbon Footprints of Intelligence}

The most immediate ethical contradiction in modern atmospheric AI is the \textbf{Energy-Climate Paradox}. To predict the impacts of global warming, we are building digital infrastructures that consume enough electricity to power small cities. 

\begin{figure}[H]
\centering
\begin{tikzpicture}[node distance=2.5cm]
    \node (sustain) [process, fill=green!20] {Sustainability};
    \node (equity) [process, fill=blue!20, right of=sustain, xshift=2cm] {Equity};
    \node (access) [process, fill=orange!20, below of=sustain] {Access};
    \node (trust) [process, fill=red!20, below of=equity] {Trust};
    
    \node (justice) [decision, below of=sustain, xshift=2.25cm, yshift=-1.25cm] {Climate Justice};
    
    \draw [arrow] (sustain) -- (justice);
    \draw [arrow] (equity) -- (justice);
    \draw [arrow] (access) -- (justice);
    \draw [arrow] (trust) -- (justice);
\end{tikzpicture}
\caption{The Ethical Dimensions of Atmospheric AI.}
\label{fig:ethics_tikz}
\end{figure}

\section{Summary}
The success of Artificial Intelligence in Earth science will be measured not by its RMSE, but by its \textbf{Human Impact}.
